\chapter{Coordinate Systems}
\label{ch:coordinates}

A common complaint from newcomers to Vim, encountering its movement commands for the first time, is that there seem to be too many ways to move the cursor, with no evident organizing principle beyond historical accretion. This chapter argues the opposite: Vim's motions are not an unstructured pile of alternatives but several distinct coordinate systems layered over the same buffer, each addressing text at a different granularity, and each internally consistent once recognized as its own system rather than judged against the others. A reader who tries to memorize forty motions as forty independent facts is working much harder than necessary; a reader who recognizes eight or nine small coordinate systems, each with two or three members, is memorizing almost nothing at all.

\section{Character Coordinates}

The finest-grained coordinate system addresses individual characters. \texttt{h} and \texttt{l} move left and right one character at a time; \texttt{0} moves to the first column of the current line and \texttt{\$} to the last; a count prefixed to either \texttt{h} or \texttt{l} moves that many characters. This is the coordinate system a reader already met in Chapter~\ref{ch:installing}, and it is deliberately the least distinctive: character-wise motion is the fallback every other coordinate system exists to make less necessary, since moving six words by pressing \texttt{l} thirty times is possible but is exactly the kind of low-level, high-repetition operation the word, sentence, and paragraph coordinate systems below are designed to replace with a single keystroke.

\section{Line Coordinates}

A second, closely related system addresses whole lines rather than characters within them. \texttt{j} and \texttt{k} move down and up one line; \texttt{G} moves to the last line of the file, and a count prefixed to \texttt{G} (\texttt{25G}) moves to that specific line number; \texttt{gg} moves to the first line. Line coordinates and character coordinates interact in a way worth noting explicitly: \texttt{j} and \texttt{k} preserve, where possible, the cursor's original column as it moves between lines of differing length, which is why moving down through a ragged block of text does not steadily drift the cursor toward column zero the way a naive implementation might.

\section{Word Coordinates}

Word motions address a coarser unit than characters, and Vim in fact defines two related but distinct notions of "word." A lowercase \texttt{w} moves to the beginning of the next word, where a word is bounded by whitespace or by a transition between alphanumeric and punctuation characters, meaning \texttt{don't} is treated as three words (\texttt{don}, \texttt{'}, \texttt{t}) under this definition. An uppercase \texttt{W} moves to the beginning of the next WORD (Vim's own capitalized term, adopted here for consistency with the documentation), where a WORD is bounded only by whitespace, treating \texttt{don't} as a single unit. \texttt{b} and \texttt{B} are the backward equivalents of \texttt{w} and \texttt{W}; \texttt{e} and \texttt{E} move to the end of the current or next word or WORD rather than its beginning; \texttt{ge} and \texttt{gE} move backward to the end of the previous word or WORD.

The word/WORD distinction is not a redundancy to be memorized as two arbitrary commands; it is a single coordinate system offered at two different resolutions, chosen depending on whether punctuation should count as a boundary. A programmer navigating through \texttt{some\_function(arg1, arg2)} will usually want the punctuation-sensitive lowercase forms, since each argument and each punctuation mark is a meaningful unit on its own; a writer navigating through ordinary prose will often find the coarser WORD forms move more efficiently through contractions and hyphenated words without stopping at every internal punctuation mark.

\section{Sentence and Paragraph Coordinates}

Moving outward in granularity again, \texttt{(} and \texttt{)} move to the beginning of the current or next sentence, and \texttt{\{} and \texttt{\}} move to the beginning of the current or next paragraph. Vim's definitions of these units are syntactic rather than semantic: a sentence, for the purposes of these motions, ends at a \texttt{.}, \texttt{!}, or \texttt{?} followed by whitespace or the end of a line (with some further trailing-punctuation handling covered in the built-in help), and a paragraph is delimited by blank lines, not by any deeper notion of topical coherence. This is worth stating plainly because it is a common source of mild disappointment: these motions cannot know that a sentence is "really" still going after an abbreviation like "Dr." followed by a capital letter, and they will treat that period as a sentence boundary regardless. The motions are fast and reliable for what they actually compute, which is punctuation and blank-line structure, and are not attempting anything more linguistically sophisticated than that.

\section{Section Coordinates}

A coarser and considerably less commonly used coordinate system addresses sections, delimited by a form-feed character or, by convention in some file types, by a line beginning with an open brace in the first column. \texttt{[[} moves backward to the beginning of the current or previous section and \texttt{]]} moves forward to the next. This system predates most modern uses of Vim and is of limited relevance to prose or typical source code, but it persists as a documented motion, and readers working with older-style C source files or documents that use form feeds as chapter separators will find it still functions exactly as originally designed.

\section{Search Coordinates}

Search introduces a coordinate system defined not by any fixed structural unit but by wherever a pattern next matches. \texttt{/pattern} searches forward, \texttt{?pattern} searches backward, and \texttt{n} and \texttt{N} repeat the most recent search in the same and opposite direction respectively. A closely related, finer-grained variant restricts the search to the current line and to a single character: \texttt{f}\textit{char} and \texttt{F}\textit{char} search forward and backward for the next occurrence of \textit{char} on the current line, landing on it, while \texttt{t}\textit{char} and \texttt{T}\textit{char} do the same but land one character before it (Chapter~\ref{ch:language} already introduced \texttt{;} and \texttt{,} as the repetition commands for this specific family). What unifies search into a single coordinate system, despite these two granularities, is that both forms address text by content rather than by a fixed unit of structure: the destination is wherever the pattern or character is found, not a position defined independently of the buffer's actual contents.

\section{Screen Coordinates}

A further coordinate system addresses position relative to the currently visible screen rather than to the buffer's absolute content. \texttt{H} moves to the first line currently visible on screen, \texttt{M} to the middle visible line, and \texttt{L} to the last visible line; each accepts a count that offsets from the relevant edge rather than jumping to an absolute buffer position. This system exists because "the top of what I am currently looking at" and "line 1 of the file" are frequently different positions, especially in a long file scrolled well past its beginning, and screen coordinates let the user address the former directly without first determining what absolute line number it happens to correspond to.

\section{Absolute Coordinates}

Distinct again from screen-relative position is absolute position within the buffer, addressed by line number directly. \texttt{nG} moves to line \textit{n}; \texttt{n\textbar} moves to column \textit{n} of the current line; \texttt{G} alone (equivalent to specifying the last line) moves to the end of the file, and \texttt{gg} (equivalent to line 1) moves to its beginning. This is the coordinate system most directly useful when working from external information, such as a compiler error message or a search tool's output, that reports a specific line number the user needs to reach directly rather than by relative navigation.

\section{Marks as a Coordinate System}

Finally, marks, already introduced in Chapter~\ref{ch:ontology} as a named entity, function grammatically as their own coordinate system once set. \verb|'a| moves to the line of mark \texttt{a}; \verb|`a| (using the backtick rather than the apostrophe) moves to the exact character position of mark \texttt{a} rather than merely its line, a distinction that matters whenever a motion or operator needs to respect the mark's precise column, not just the line it sits on. Two automatically maintained marks are worth knowing from the outset: \verb|``| (or, equivalently, \verb|''|) returns to the position before the most recent jump, and \verb|'.| returns to the position of the most recent change, both covered further in Part IV's chapter on Marks and the jump list. What makes marks a coordinate system rather than simply a storage mechanism is precisely this: once set, a mark is addressable by every operator exactly the way a word or a line number is, filling the same grammatical slot established in Chapter~\ref{ch:language}.

\section{Choosing a Coordinate System}

The practical skill this chapter is building toward is not memorizing every motion in isolation but developing the reflex to ask, given an editing goal, which coordinate system most directly addresses it. Deleting to the end of a function call is a job for word or bracket-matching coordinates (the latter covered alongside structural regions in Chapter~\ref{ch:regions}); deleting to a specific recalled position elsewhere in the file is a job for a mark; deleting to the next occurrence of a particular character on the current line is a job for \texttt{f} or \texttt{t}. Fluency in Vim looks, from the outside, like choosing the right motion instantly and without visible deliberation; from the inside, it is closer to routing an editing goal to the coordinate system that already matches its shape, a skill this chapter's organization is intended to make explicit rather than leaving it to be inferred gradually from years of accumulated habit.
